\subsection{Rank of Smooth Maps}

given a smooth $f:U\gto V$ for $U\subseteq\bR^{n}$ and $V\subseteq\bR^{m}$,
the \textbf{rank} of $f$ at $p$ is the rank of $D_{p}f$.
we prefer to think of the jacobian as a linear map from $\bR^{n}$ to $\bR^{m}$.
we can equivalently defn the derivative of $f$ at $p$ as:
\begin{equation*}
	D_{p}f:\bR^{n}\gto\bR^{m} \qquad v \ \mto \ \frac{\d}{\d t}\bigg\rvert_{t=0}f(p+tv)
\end{equation*}
where $t\in\bR$. the rank of $f$ at $p$ is then the rank of this linear map.
it may be sometimes helpful to recall that diffeos are full rank.
we now generalize.

\begin{defn}
	let $f\in C^{\infty}(M,N)$, $p\in M$. the \textbf{rank} of $f$ at $p$ is then:
	\begin{equation*}
		\trm{rank}_{p}(f) \ = \ \trm{rank}_{\vphi(p)}(\psi\circ f\circ\vphi^{-1})
	\end{equation*}
	for appropriate charts. the point $p$ is a \textbf{regular pt} if $f$ has
	maximal rank at $p$; it is a \textbf{critical pt} otw.
\end{defn}
in other words, we simply apply local coords then use our usual defn.
furthermore, this is clearly independent of choice of chart, since once again,
transition maps are diffeos.

we now examine the rank of maps $C^{\infty}(M,N)$.
we separate into 3 cases based on the dims of $M$ and $N$.

first, we consider when $m=n$. we aim to generalize the inverse fn theorem.
\begin{thm}[title=Inverse Function Theorem]
	let $f\in C^{\infty}(M,N)$ where $m=n$. if $r_{p}(f)=m$ for some $p$, then
	there exists open nbhd $U\subseteq M$ of $p$ such that $f$ restricts to a diffeo on $U$.
\end{thm}
proof: apply regular inv fn thm. oh btw $r_{p}(f)$ is rank of $f$ at $p$.

\begin{defn}
	$f\in C^{\infty}(M,N)$ is called a \textbf{local diffeo} if $\fa p\in M$,
	$f$ restricts to a diffeo on some open nbhd $U$ of $p$.
	equivalently, $f$ is full rank everywhere.
\end{defn}
as an example, the quot map $\pi:S^{n}\gto\bR P^{n}$ is a local diffeo on hemispheres.

next, we consider when $m\geq n$. we aim to generalize the implicit fn thm.
\begin{thm}[title=this has many names]
	let $f\in C^{\infty}(M,N)$ and $p\in M$ s.t. $r_{p}(f)=n$.
	then $\ex$ appropriate charts s.t.:
	\begin{equation*}
		(\psi\circ f\circ\vphi^{-1})(x_{1},x_{2})=x_{1}
	\end{equation*}
	for all $x=(x_{1},x_{2})$ in $\vphi(U)$.
	in particular, $f^{-1}(q)\cap U$ is a submfld of codim $n$.
\end{thm}
pf: apply implicit fn thm. supposedly. i still dk what that thm is.

\begin{defn}
	a pt $q\in N$ is a \textbf{reg value} of $f$ if $r_{p}(f)=n$ for all $p\in f^{-1}(q)$.
	it is a \textbf{critical value} or \textbf{singular value} otw.
	note that pts not in the img are considered reg vals.
\end{defn}
we can then restate 6.4 as:
\begin{thm}
	for any reg val $q\in N$ of $f$, $f^{-1}(q)$ is a submfld of codim $n$.
\end{thm} \

\begin{defn}
	$f$ is a \textbf{submersion} if it is full rank everywhere.
\end{defn}
as an ex, proj maps $\pi:M\times N\gto M$ are submersions; 6.4 says that locally,
\textit{every} submersion looks like this.

finally, consider the case where $m\leq n$.
what we aim to generalize is not named lol.
\begin{thm}[title=Normal Form for Immersions (supposedly)]
	let $f\in C^{\infty}(M,N)$ and $p\in M$ s.t. $r_{p}(f)=m$.
	then $\ex$ appropriate charts s.t.:
	\begin{equation*}
		(\psi\circ f\circ\vphi^{-1})(x)=(x,0)
	\end{equation*}
	for $x\in\vphi(U)$. in particular $f(U)\subseteq N$ is a submfld of dim $m$.
\end{thm}
pf is same idea as always. again, never heard of this thm even in $\bR^{n}$ tho.

\begin{defn}
	$f$ is an \textbf{immersion} if $r_{p}(f)=m$ for all $p\in M$.
\end{defn}
for instance, (smooth) curves are immersions.
the book calls embeddings immersions ``given as the inclusion map for a submfld".
more commonly, it is an inj immersion which is topologically an embedding, i.e. diffeo onto its img.

\begin{thm}
	sps $f:M\gto N$ an inj immersion, $M$ cpt.
	then $f(M)\subseteq N$ an embedded submfld.
\end{thm} \

\begin{pf}[source=Primary Source Material]
	we want to find submfld charts.
	take the chart in $N$ obtained from thm 6.8, and shrink it to be disjoint
	from $f(M\sm U)$.
\end{pf}

to summarize this section, we showed that if a smooth map $f$ has
maximal rank near a pt $p$, we can choose local coords arnd $p$ and $f(p)$
s.t. $f$ in those coords becomes linear.
in particular, near any pt of $M$, submersions look like surj lin. maps, and
immersions look like inj lin. maps.
